\section{流模型的优势与推广能力}

在建立了等价关系之后，讲者转向讨论流模型相对于扩散模型的\textbf{独特优势}——即流模型在哪些方面真正``超越''了扩散模型。

\subsection{非线性条件流映射}

扩散模型本质上使用的是线性高斯形式的条件路径。而 Flow Matching 框架允许定义\textbf{任意}的条件流映射，不限于线性形式：

\begin{itemize}
    \item 可以尝试非线性的 $\psi_t(\mathbf{x}_0 \mid \mathbf{x}_1)$，例如加入余弦、sigmoid 等非线性变换
    \item 实验表明，虽然非线性流映射目前并没有显著超越线性形式，但框架本身的灵活性为未来探索提供了空间
\end{itemize}

\subsection{任意参考分布}

这是流模型最重要的推广之一。扩散模型要求参考分布必须是标准高斯分布 $\mathcal{N}(\mathbf{0}, \mathbf{I})$，因为整个理论（前向扩散过程）建立在高斯噪声的基础上。而流模型可以使用\textbf{任意}参考分布：

\begin{importantbox}{分布到分布的映射}
Flow Matching 允许学习从\textbf{任意}源分布到\textbf{任意}目标分布的映射。这不仅包括：
\begin{itemize}
    \item 标准高斯 $\to$ 数据分布（经典生成任务）
\end{itemize}
还包括：
\begin{itemize}
    \item 数据分布 A $\to$ 数据分布 B（风格迁移、域转换）
    \item 低分辨率分布 $\to$ 高分辨率分布（超分辨率）
    \item 任意复杂分布之间的转换
\end{itemize}
这种灵活性是扩散模型框架所不具备的。
\end{importantbox}

\subsection{ODE vs.\ SDE 的出发点}

\begin{knowledgebox}{两种建模哲学}
扩散模型和流模型的另一个根本区别在于建模的\textbf{出发点}：
\begin{itemize}
    \item \textbf{扩散模型}：从\textbf{随机微分方程}（SDE）出发，将前向扩散定义为逐步加高斯噪声的随机过程；然后发现等价的\textbf{确定性} ODE 也能给出相同的边际分布
    \item \textbf{流模型}：直接从\textbf{常微分方程}（ODE）出发，定义从源分布到目标分布的确定性映射，无需引入随机性
\end{itemize}
流模型的出发点更简洁，避免了 SDE 到 ODE 转换的复杂推导。
\end{knowledgebox}

\subsection{本章小结}

流模型在以下方面推广了扩散模型：(1) 条件流映射不限于线性形式；(2) 参考分布不限于标准高斯；(3) 直接从 ODE 出发，概念更简洁。这些推广为更灵活的生成模型设计提供了理论基础。

